\documentclass[10pt]{article}

\usepackage[utf8]{inputenc}
\usepackage[T1]{fontenc}
\usepackage[english]{babel}
\usepackage[left=1.25cm,top=0.90cm,right=1.25cm,bottom=0.40cm]{geometry}
\usepackage{enumitem}
\usepackage{xcolor}
\usepackage{hyperref}
\usepackage{titlesec}
\usepackage{tabularx}
\usepackage{array}
\usepackage{microtype}
\usepackage{lmodern}

\setlength{\parindent}{0pt}
\setlength{\parskip}{0pt}
\linespread{1.105}
\pagestyle{empty}

\definecolor{accent}{HTML}{1F4E79}
\definecolor{muted}{HTML}{596675}
\definecolor{rulegray}{HTML}{CAD3DC}

\hypersetup{
    colorlinks=true,
    urlcolor=accent,
    linkcolor=accent,
    pdfauthor={Nguyen Minh Khoa},
    pdftitle={Nguyen Minh Khoa - AI Engineer and Researcher}
}

\newcommand{\cvsection}[1]{%
    \vspace{6.2pt}%
    {\color{accent}\bfseries\small #1}\par
    \vspace{2.3pt}%
    {\color{rulegray}\hrule height 0.45pt}\vspace{4.0pt}%
}

\newcommand{\entry}[2]{%
    \textbf{#1}\hfill {\color{muted}\small #2}\par
}

\newcommand{\subentry}[1]{%
    {\color{muted}\small #1}\par
}

\setlist[itemize]{
    leftmargin=13pt,
    labelsep=5pt,
    itemsep=2.8pt,
    topsep=2.8pt,
    parsep=0pt,
    partopsep=0pt
}

\begin{document}

\begin{center}
    {\LARGE\textbf{NGUYEN MINH KHOA}}\\[-1pt]
    {\color{accent}\textbf{AI Engineer \& Researcher}}\\[3pt]
    {\small Ho Chi Minh City, Vietnam \textbullet\
    \href{tel:+84907713819}{+84 907 713 819} \textbullet\
    \href{mailto:ngminhkhoayj2706@gmail.com}{ngminhkhoayj2706@gmail.com} \textbullet\
    \href{https://www.linkedin.com/in/ngminhkhoa1410}{LinkedIn} \textbullet\
    \href{https://github.com/Purin1410}{GitHub} \textbullet\
    \href{https://scholar.google.com/citations?user=tCq7yoQAAAAJ}{Google Scholar} \textbullet\
    \href{https://purin1410.github.io}{Portfolio}}
\end{center}

\vspace{-2pt}

\cvsection{PROFILE}
Applied AI engineer and Research Assistant at AiTA Lab, FPT University, bringing model development, controlled experimentation and application engineering to AI product development. Experience spans text-to-molecule generation, mathematical expression recognition and LLM research workflows, connecting training data and evaluation with deployed inference services.

\cvsection{RESEARCH EXPERIENCE}
\entry{Research Assistant | AiTA Lab, FPT University}{Oct 2024 - Present}
\begin{itemize}
    \item Developed and evaluated deep-learning methods across HMER and text-to-molecule generation, covering data preparation, model training, ablation studies, and reproducible experiment tracking.
    \item Selected published outcomes include up to \textbf{+5 ExpRate points} for Mask CoMER; ChemAligner-T5 reports \textbf{69.77\% BLEU} and \textbf{31.28\% Levenshtein distance} on L+M-24.
\end{itemize}

\cvsection{SELECTED PROJECTS}
\entry{LexiChem | AI \& Inference Contributor}{2025 - 2026}
\subentry{Deployed molecular generation web application with FastAPI and NVIDIA Triton model serving.}
\begin{itemize}
    \item Built a three-dataset preprocessing pipeline covering \textbf{693,090 training examples} across L+M-24, Mol-Instructions, and ChEBI-20, with RDKit validation and canonicalization, deduplication, and SMILES-to-SELFIES conversion.
    \item Capstone ChEBI-20 experiments with shared-latent alignment improved BioT5+ from \textbf{87.2 to 88.8 BLEU}, \textbf{52.2 to 52.9 exact match}, and \textbf{90.7 to 93.5 MACCS FTS}, while maintaining \textbf{100\% validity}.
    \item Integrated multi-model inference with FastAPI, NVIDIA Triton, RDKit, and MongoDB for molecule generation, validation, comparison, and saved experiment runs.
\end{itemize}

\entry{Research Ops}{2026 - Present}
\subentry{Research infrastructure orchestrating literature understanding, experiment execution and evidence-backed manuscript development.}
\begin{itemize}
    \item Architected a FastAPI/PostgreSQL execution engine for configurable LLM research workflows, with persistent scientific state, execution budgets and human decision gates.
    \item Linked grounded LLM execution, experiment artifacts and manuscript evidence through versioned records and recovery.
\end{itemize}

\entry{NexOps | AI, Simulation \& Platform Lead}{2025 - Present}
\begin{itemize}
    \item Built a MES/SCADA prototype with repeatable machine/shift simulations and MQTT/EMQX telemetry through FastAPI and TimescaleDB for production planning, live monitoring, alarms and OEE.
    \item Integrated Bambu Lab printer control over LAN; coordinated hardware messaging with an IoT teammate.
\end{itemize}

\cvsection{SELECTED PUBLICATIONS}
\textbf{LexiChem: Shared-Latent Alignment for Faithful Text-to-Molecule Generation in SELFIES Space}\\[-1pt]
{\small APWeb-WAIM 2026 | First author | Accepted and presented; publication pending}\par
\vspace{2.2pt}
\textbf{Mask CoMER: Enhancing Handwritten Mathematical Expression Recognition with Masked Language Pretraining and Regularization}\\[-1pt]
{\small ICDAR 2025 | \textbf{ICORE 2026 Rank A} | Co-author | Published}\par
\vspace{2.2pt}
\textbf{Enhancing handwritten mathematical expression recognition with convolutional block attention refinement module and multi-task learning}\\[-1pt]
{\small Journal of Information and Telecommunication | \textbf{SCImago Q2 (2025)} | Published online Aug 2026}

\cvsection{EDUCATION}
\entry{Bachelor of Science in Artificial Intelligence | FPT University, HCMC}{2022 - 2026}

\cvsection{AWARDS \& SCHOLARSHIPS}
{\small \textbf{50\% Tuition Scholarship}, FPT University (2022) \textbullet\
\textbf{RESFES HCMC Second Prize} (2025, 2026) \textbullet\
\textbf{Demo Pitching Day} (VND 50M team award/project funding, 2025) \textbullet\
\textbf{Innovation Quest} (VND 40M team award/funding, 2025)}

\cvsection{TECHNICAL SKILLS}
{\footnotesize
\textbf{AI \& Research:} Python, PyTorch, PyTorch Lightning, Hugging Face Transformers, scikit-learn, OpenCV, RDKit, ablation studies, error analysis, Weights \& Biases, MLflow\\
\textbf{Backend \& Data:} FastAPI, Pydantic, Pandas, PostgreSQL, TimescaleDB, MongoDB, Redis, MQTT/EMQX, REST APIs, NVIDIA Triton\\
\textbf{Engineering \& Reproducibility:} Docker, Docker Compose, Linux, Git, Nginx, Playwright, LaTeX, scientific writing
}

\end{document}
